\documentclass[10pt,a4paper]{amsart}
\usepackage[margin=19mm]{geometry}
\usepackage{amsmath,amssymb,mathtools}
\usepackage[scr=rsfs,cal=euler]{mathalfa}
\usepackage{xfrac}
\usepackage[unicode,hidelinks,bookmarks=false]{hyperref}
\newcommand{\ie}{i.e.\ }
\newcommand{\cf}{cf.\ }
\newcommand{\resp}{respectively\ }
\newcommand{\Subsection}{\subsection}
\providecommand{\nobreakdash}{\nobreak\hbox{-}}
\setlength{\emergencystretch}{2em}
\multlinegap0pt
\usepackage{amssymb}
\usepackage{tikz}
\usepackage{algorithm}\usepackage{algpseudocode}
\usepackage{enumerate}
\usepackage{mathtools}
\newtheorem{lemma}{Lemma}
\newtheorem{theorem}{Theorem}
\newtheorem{claim}{Claim}
\title{t-tone colorings of outerplanar and Halin graphs}
\author{Hadeel Al Bazzal, Olivier Togni}
\date{Source: arXiv:2603.18674v1 (CC BY 4.0)}
\begin{document}
\maketitle

\noindent\emph{Source and licence.} t-tone colorings of outerplanar and Halin graphs. Version arXiv:2603.18674v1 (CC BY 4.0), distributed by arXiv under the Creative Commons Attribution 4.0 International licence, http://creativecommons.org/licenses/by/4.0/, as recorded in the arXiv licence metadata for this exact version and shown on the abstract page https://arxiv.org/abs/2603.18674v1. Full licence notice: \texttt{licenses/arxiv-2603.18674-CC-BY-4.0.txt}.\par\smallskip

\noindent\textit{Editorial scope.} Counted unit: the article's theorem theorem (source0.tex lines 288--290). The article's complete original proof of it is delivered with it (source0.tex lines 292--336, one \texttt{proof} environment of 45 lines). No mathematics is written, repaired or re-derived: the statement and the proof are the article's own lines, in the article's own notation, and no retained line is substituted or re-flowed.  Retained uncounted as the source-local context the proof needs: lines 280--282.  Nothing was omitted from any retained span, so the package carries no omission record.  No retained line contains a citation, so the wrapper carries no bibliography and no bibliography entry.  No retained line contains a figure environment, an image-inclusion command or drawing code, so no image is delivered and the package declares no asset dependency.  Editorial formatting only, in addition: an AMS \texttt{amsart} preamble that restates the article's own declarations and macros used by the retained lines, namely the environments \texttt{lemma}, \texttt{theorem}, \texttt{claim} on separate counters, together with the standard packages those lines need.  The retained environments print in this document as Lemma 1, Theorem 1, Claim 1 in order of appearance, whereas the article prints the same items with the numbering of its own full document.  The complete native source of the article as distributed in the arXiv e-print for this exact version is bundled unchanged as \texttt{sources/196711/source0.tex}.
\begin{lemma}\label{l3}
Every cycle admits a $3$-good $11$-coloring.
\end{lemma}

\begin{theorem}
    Every subcubic outerplanar graph $G$ has a $3$-good $11$-coloring.
\end{theorem}

\begin{proof} 
Suppose, for the sake of contradiction, that the statement is false, and let $G$ be a counterexample of minimum order. Clearly, $G$ is connected.
 
\begin{claim}
    $\delta(G) \ge 2$.
\end{claim}

\begin{proof}
Assume for contradiction that $G$ contains a vertex $u$ of degree one, and let $v$ be its single neighbor. Consider the graph $G' = G - \{u\}$. By the minimality of $G$, the smaller graph $G'$ admits a $3$-good $11$-coloring, denoted by $f$.\\
\noindent If $v$ has only one neighbor $v'$ in $G'$, then choose one color from the label of $v'$ and select two from the remaining colors from $\{1,\dots,11\}$ not used by $f(v)$ and $f(v')$, and assign these three colors to $u$. This yields a valid extension of $f$, a contradiction.\\
\noindent Suppose instead that $v$ has two distinct neighbors, $v'$ and $v''$. Since the distance between $v'$ and $v''$ is at most two, their two labels are distinct. We may then assign to $u$ three colors: $a \in f(v')$, $b \in f(v'')$ with $a \ne b$, and $c \in \{1,\dots,11\}\setminus \bigl(f(v)\cup f(v')\cup f(v'')\bigr)$. This extends the coloring to a $3$-good $11$-coloring of $G$, yielding a contradiction.\end{proof}

\noindent Note that $G$ must contain at least two bounded faces; otherwise, $G$ would be a cycle, which admits a $3$-good $11$-coloring by Lemma \ref{l3}, a contradiction. Let $F$ be a pendant face in $G$ with $C(F) = v_1v_2 \dots v_\ell v_1$, where $\ell \ge 3$. By the definition of a pendant face in a subcubic outerplanar graph, we assume that $v_1$ and possibly $v_\ell$ are the only vertices of degree $3$, with all others having degree $2$.

\begin{claim}
   $\ell \leq 4$.
\end{claim}

\begin{proof}
Suppose, for a contradiction, that $\ell \ge 5$. Consider the graph 
$G' = (G - \{v_3\}) + v_2v_4$. By the minimality of $G$, the graph $G'$ admits a $3$-good $11$-coloring $f$. Without loss of generality, assume that $f(v_2) = 123$ and $f(v_4) = 456$, and let $S = \{7,8,9,10,11\}$. Since $f$ is a $3$-good $11$-coloring of $G'$, the vertex $v_5$ shares exactly one color with $v_2$ and takes its two remaining colors from $S$. Thus, we may assume $f(v_5) = 178$. Similarly, $v_1$ shares exactly one color with $v_4$, say $4$, and takes its two remaining colors from $S$. Hence, the sets $f(v_1), f(v_2), f(v_4)$, and $f(v_5)$ together use at most ten colors, leaving one unused color from $\{1,\dots,11\}$, which we may assume to be $11$. Observe that $v_1$ may share two, one, or no colors with $v_5$. Accordingly, we may assume that 
$f(v_1) = 478$, $479$, or $4910$, respectively. In all cases, we assign $f(v_3) = 8910$.

\noindent It remains to ensure that $v_2$ and $v_4$ share a color in $G$, which we achieve by relabeling $v_4$. Note that $d_G(v_5) \ge 2$.

\noindent Suppose first that $d_G(v_5)=2$, and let $x$ be a neighbor of $v_5$ distinct from $v_4$. In this case, $x \in \{v_1, v_6\}$, and $x$ together with $v_2$ are the only second neighbors of $v_4$ in $G$. Since $f$ is a $3$-good $11$-coloring, $v_4$ and $x$ share exactly one color. Without loss of generality, assume that this color is $4$. If $2 \in f(x)$, we relabel $f(v_4)$ as $256$; otherwise, we relabel it as $245$. In either case, the coloring extends to $G$, yielding a contradiction.

\noindent Thus, $d_G(v_5)=3$. It follows that $\ell=5$, and hence $v_1$ and $v_5$ are adjacent, which fixes $f(v_1)=4910$. Let $x$ be the third neighbor of $v_5$ in $G$, distinct from $v_1$ and $v_4$. In this case, $v_4$ has exactly three second neighbors in $G$, which are $v_1$, $v_2$, and $x$. Observe that $v_1$ and $v_4$ share the color $4$. Moreover, $x$ and $v_4$ share exactly one common color; denote it by $a$.

\noindent Note that $v_1$ and $x$ are not adjacent. Otherwise, $x$ and $v_2$ would be at distance two in $G$, and hence must share exactly one common color other than $1$. Without loss of generality, assume this color is $2$. Then $3 \notin f(x)$. Moreover, $a \neq 4$; assume without loss of generality that $a=5$. Now relabel $v_4$ with $345$, which extends $f$ to $G$, a contradiction.

\noindent Thus, $x$ and $v_1$ are at distance two in $G$, and hence must share exactly one common color. First, suppose that this common color is $4$. If $2 \in f(x)$, we relabel $f(v_4)$ as $2510$; otherwise, we relabel it as $245$. In either case, $f$ extends to $G$, a contradiction.

\noindent Therefore, we may assume, without loss of generality, that $x$ shares the color $5$ with $v_4$ and the color $9$ with $v_1$. Consequently, $|\{2,3\} \cap f(x)| \leq 1$. Suppose, without loss of generality, that $2 \notin f(x)$. We can then relabel $f(v_4)$ as $245$, which once again leads to a contradiction.\end{proof}

\begin{claim}
    $\ell= 3$
\end{claim}

\begin{proof}
Suppose, to the contrary, that $\ell= 4$; then let $G' = (G - \{v_2\}) + v_1v_3$. Without loss of generality, suppose $f(v_1) = 123$ and $f(v_3) = 456$. Set, without loss of generality, $f(v_4) = 789$. Let $y$ be the neighbor of $v_1$ in $G$ other than $v_2$ and $v_4$. Since $f$ is a $3$-good $11$-coloring, $y$ shares exactly one color with $v_3$. Suppose $v_4$ and $y$ are not adjacent; then $y$ also shares exactly one color with $v_4$. Hence, without loss of generality, let $f(y) = 4710$. If $d_G(v_4)=2$, then relabel $v_3$ as $156$ and label $v_2$ as $81011$, a contradiction. Thus, let $y'$ be the third neighbor of $v_4$. Since $G'$ admits a $3$-good $11$-coloring, $y'$ shares one color with $v_1$, say $1$, and one with $v_3$, say $a$. Relabel $f(v_3)$ by replacing one of its colors, distinct from $a$, with $2$, and assign $f(v_2)=81011$, yielding a contradiction. Thus, $v_4$ and $y$ are adjacent. Therefore, without loss of generality, let $f(y) = 41011$. Hence, we can relabel $v_3$ as $426$ and then label $v_2$ as $5710$, a contradiction.\end{proof}

\noindent Then $F$ is a triangle. Suppose first, without loss of generality, that $d_G(v_1)=3$ and $d_G(v_2)=d_G(v_3)=2$. Let $G' = G - \{v_2, v_3\}$. By the minimality of $G$, the graph $G'$ admits a $3$-good $11$-coloring $f$. Let $y$ denote the neighbor of $v_1$ in $G'$. Suppose, without loss of generality, that $f(v_1)=123$ and $f(y)=456$. Assigning the labels $478$ and $5910$ to $v_2$ and $v_3$, respectively, yields a $3$-good $11$-coloring of $G$, a contradiction.

\noindent Therefore, we may assume, without loss of generality, that $d_G(v_1)=d_G(v_2)=3$ and $d_G(v_3)=2$. Let $G' = G - \{v_3\}$. By the minimality of $G$, the graph $G'$ admits a $3$-good $11$-coloring $f$. Suppose, without loss of generality, that $f(v_1)=123$ and $f(v_2)=456$. Let $y_1$ and $y_2$ be the neighbors of $v_1$ and $v_2$, not in $C(F)$, respectively. If $y_1=y_2$, then we can set $f(y_1)=789$. Now, assign the label $71011$ to $v_3$, a contradiction. Thus, $y_1$ and $y_2$ are distinct. Let $S=\{7,8,9,10,11\}$. Because $G'$ admits a $3$-good $11$-coloring, $y_2$ shares exactly one color with $v_1$ and takes its two remaining colors from $S$. Hence, without loss of generality, we may set $f(y_2) = 178$. Similarly, $y_1$ shares exactly one color with $v_2$, say $4$, and takes its two remaining colors from $S$. Therefore, the four sets $f(v_1), f(v_2), f(y_1), f(y_2)$ together involve at most ten colors, leaving one unused in $\{1,\dots,11\}$, which we may assume, without loss of generality, to be $11$. Observe that $y_1$ may share two, one, or no colors with $y_2$. Without loss of generality, we may assume that $f(y_1) = 478$, $479$, or $4910$, respectively. In all cases, we assign $f(v_3) = 8911$ and obtain a $3$-good $11$-coloring of $G$, a contradiction.\end{proof}

\end{document}
